\documentclass[book.tex]{subfiles}

\begin{document}

\chapter{Schrodinger's Equation in Three Dimensions}

\label{chap:schrodinger_3D}
\noindent\rule{11cm}{0.4pt} \\
\textbf{Prerequisites:}  \autoref{chap:schrodinger} \\
\textbf{Difficulty Level:} ** \\
\noindent\rule{11cm}{0.4pt}
\vspace{5mm}

In this chapter, we study techniques for solving Schrodinger's equation in three dimensions. We can write the equation as

\begin{align}
-\frac{\hbar^2}{2m}\left [ \frac{\partial}{\partial x^2} + \frac{\partial}{\partial y^2} + \frac{\partial}{\partial z^2}\right ] \psi(x, y, z) + V(x, y, z) \psi(x, y,z) &= E \psi(x, y, z) \\
-\frac{\hbar^2}{2m} \nabla^2 \psi + V \psi &= E \psi
\end{align}

\section{Moving to Spherical Coordinates}
% See https://math.stackexchange.com/questions/3734106/laplacian-of-spherical-coordinates
The three dimensional Schrodinger's equation becomes easier to solve in spherical coordinates. A technical challenge to performing this operation is that we need to represent the Laplacian operator $\nabla^2$ in spherical coordinates. Let us introduce the change of variables

%\textcolor{red}{TODO: Do the derivation for general curvilinear coordinates and then specialize. The direct derivation is very painful.}

\begin{align}
    u &= u(x, y, z) \\
    v &= v(x, y, z) \\
    w &= w(x, y, z)
\end{align}
In reverse we can write
\begin{align}
    x &= x(u, v, w) \\
    y &= y(u, v, w) \\
    z &= z(u, v, w)
\end{align}

Let 
\begin{align}
    r = (x, y, z)
\end{align}
Then we write
\begin{align}
    dr &= (dx, dy, dz)
\end{align}
We can also write
\begin{align}
    dr &= \frac{\partial r}{\partial u} du + \frac{\partial r}{\partial v} dv + \frac{\partial r}{\partial w} dw
\end{align}

We write
\begin{align}
    h_u &= \left | \frac{\partial r}{\partial u} \right | \\
    h_v &= \left | \frac{\partial r}{\partial v} \right | \\
    h_w &= \left | \frac{\partial r}{\partial w} \right | 
\end{align}

These terms are typically referred to as scale factors. The gradient term in the $(u, v, w)$ coordinate system is then given by

\begin{align}
\nabla &= \left ( \frac{1}{h_u} \frac{\partial}{\partial u}, \frac{1}{h_v} \frac{\partial}{\partial v}, \frac{1}{h_w} \frac{\partial}{\partial w} \right )
\end{align}
For the specific case of spherical coordinates, we have a coordinate system $(r, \theta, \phi)$ defined by the relationship
\begin{align}
x &= r \cos \theta \sin \phi
\end{align}
We can then compute that the gradient operator is given by
\begin{align}
\nabla &= \left ( \frac{\partial}{\partial r}, \frac{1}{r \sin \phi} \frac{\partial}{\partial \theta}, \frac{1}{r} \frac{\partial}{\partial \phi} \right) 
\end{align}

Computing the divergence operator in curvilinear coordinates is more difficult. The general formula which we do not prove is given by
\begin{align}
\nabla \cdot f &= \left ( h_u h_v h_w \right ) \left ( \frac{\partial}{\partial u} \left ( h_v h_w  x \right) + \frac{\partial}{\partial v} \left ( h_u h_w  y \right) + \frac{\partial}{\partial w} \left ( h_u h_v  z \right) \right )
\end{align}

Combining this formula with the previous formula allows for the derivation of the spherical Laplacian equation we use in the next section.
%\textcolor{red}{TODO: Compute the gradient and the divergence in general curvilinear coordinates and obtain Laplacian. Then substitute in the spherical coordinates. Derivation notes are in the google drive.}

\section{Solving the Equation}

For simplicity, we assume that the system we are modeling is time independent. In this case, Schrodinger's equation takes the form 

\begin{align}
-\frac{\hbar^2}{2m} \nabla^2 \psi + V \psi = E \psi
\end{align}

We use the spherical form of the Laplacian %\textcolor{red}{TODO: Add a derivation of the spherical Laplacian}.

\begin{align}
    \nabla^2 = \frac{1}{r^2} \frac{\partial}{\partial r} \left ( r^2 \frac{\partial}{\partial r} \right) + \frac{1}{r^2 \sin \theta} \frac{\partial}{\partial \theta} \left ( \sin \theta \frac{\partial}{\partial \theta} \right ) + \frac{1}{r^2 \sin^2 \theta} \left ( \frac{\partial^2}{\partial \phi^2}\right )
\end{align}

Plugging this into Schrodinger's equation, we obtain

\begin{align}
    -\frac{\hbar^2}{2m} \left [ \frac{1}{r^2} \frac{\partial}{\partial r} \left ( r^2 \frac{\partial \psi}{\partial r} \right) + \frac{1}{r^2 \sin \theta} \frac{\partial }{\partial \theta} \left ( \sin \theta \frac{\partial \psi}{\partial \theta} \right ) + \frac{1}{r^2 \sin^2 \theta} \left ( \frac{\partial^2 \psi}{\partial \phi^2}\right ) \right ] + V \psi = E\psi
\end{align}

$\psi$ is now a function of the three variables $(r, \theta, \psi)$ and Schrodinger's equation becomes a partial differential equation in these three variables. To solve this equation we look to separation of variables. In particular, we hypothesize the following form for solutions

\begin{align}
    \psi(r, \theta, \phi) &= R(r)Y(\theta, \phi)
\end{align}

Here $R(r)$ is the radial component of the solution and $Y(\theta,\phi)$ is the angular component. Using this ansatz, Schrodinger's equation reduces to solving two separate ordinary differential equations.

%\textcolor{red}{TODO: Finish this section. SHould be a fairly standard derivation of Schrodinger's in 3D}
\end{document}